% NOTE: This Kazakh translation is a working draft. A native Kazakh-speaking
% reviewer (preferably with a technical / cybersecurity background) should
% verify terminology before submission.

\begin{kazakh}
\thispagestyle{empty}
\begin{center}
{\Large \textbf{АҢДАТПА}\par}
\end{center}
\vspace{0.5cm}

Қазіргі өндірістік басқару жүйелеріндегі (ICS) аномалияларды анықтау детекторлары өздері оқытылған деректер жиынында жоғары F1 көрсеткіштерін көрсетеді, алайда уақытқа сезімтал бағалау метрикалары қолданылғанда олардың рейтингтері айтарлықтай өзгереді, ал тест стенділері арасындағы тасымалдау кезінде олар жалпылау қабілетін жоғалтады --- бұл екі әсер де әдебиетте басым тақырыптық сандардың арасында жасырын қалады.  Бұл дипломдық жұмыс алты детекторды (Isolation Forest, One-Class SVM, тығыз автокодер, LSTM-автокодер, USAD, TranAD) HAI 21.03 және Morris (магистральдік газ құбыры) деректер жиындарында үш уақытқа сезімтал метриканы қолдана отырып салыстырады, өкілдіктің тасымалдануы мен шекті мәннің тасымалдануын ажырататын стенділер арасындағы тасымалдау зерттеуін жүргізеді, сондай-ақ технологиялық процесс деңгейінде сенсорлық атрибуцияны бағалайды.

Жұмыстың үш негізгі тұжырымы: (1)~бастапқы доменде калибрленген шекті мәнмен жүргізілген стенділер аралық тасымалдау әрбір детекторды мақсатты доменнің класс ықтималдығына сәйкес «бәрі --- шабуыл» тривиалды классификаторына дейін төмендетеді, ал мақсатты доменде калибрленгенде классикалық және заманауи модельдер арасындағы стандартты хаттамада көрінбейтін мәнді бөліну анықталады; (2)~HAI ішіндегі «бір процесті алып тастау» эксперименті белгі үлестірімінің ығысуын класс ықтималдығының ығысуынан ажыратып, процеске тән сенсорлар алынғанда мақсатты «соқырлықты» көрсетеді; (3)~тығыз автокодер --- F1 бойынша ең әлсіз детектор --- HAI-дың үш шабуылданған процесінің барлығында да белгі бойынша атрибуциясы кездейсоқ деңгейден тұрақты түрде жоғары болатын жалғыз модель, бұл анықтау мен локализация арасындағы өзара ымыраны айқындайды.

Жұмыстың ғылыми үлесі: қайталанатын салыстырмалы тестілеу хаттамасы, біріктірілген детектор интерфейсі шеңберіндегі ашық бастапқы кодты іске асыру, сондай-ақ ICS аномалияларды анықтау қауымдастығына әдіснамалық ұсыныс: бір тақырыптық F1 мәнінің орнына метрикалардың, калибрлеу режимдерінің және атрибуция нәтижелерінің толық кестесін ұсыну.  Құжаттағы әрбір сан репозиторийде сақталған YAML конфигурацияларынан қайталанады.

\vspace{0.5cm}
\noindent\textbf{Кілттік сөздер:} өндірістік басқару жүйелері; аномалияларды анықтау; кибер-физикалық қауіпсіздік; стенділер аралық тасымалдау; уақытқа сезімтал метрикалар; атрибуция; HAI; Morris.
\end{kazakh}
\newpage
